% ==========================================================
\section{Project Scope and Method}

\subsection{Statement of the Problem}
The goal of this project is to design and optimize a CUDA-based license plate detection and character recognition pipeline from an existing OpenCV-based implementation. 
The main challenge is not only kernel speed, but also maintaining high GPU utilization when processing many images, where host-device transfer latency can reduce the throughput due to transfers on the PCI-e or NVLinks.

The project is organized in phases. It starts from a CPU baseline, then migrates the full detection pipeline to CUDA, introduces dynamic batch sizing, and finally evaluates asynchronous execution and classifier behavior.

\subsection{Fixed Benchmark Protocol}
To ensure fair comparisons, all devices process the same fixed image set in the same deterministic order. The benchmark set is selected once from the dataset and reused unchanged across runs.

The following metrics are recorded per image and per batch:
\begin{itemize}
    % \item End-to-end latency (ms/image and ms/batch)
    \item Throughput (images/s)
    \item Stage-level timing (preprocess, plate detection, character detection, KNN)
    \item Transfer/compute split (H2D, kernel time, D2H)
    \item Detection output statistics (plates detected, characters recognized)
\end{itemize}

\subsection{Target Platforms}
The comparison includes one CPU baseline and three GPU targets:
\begin{itemize}
    \item CPU implementation (reference baseline)
    \item NVIDIA RTX 4070 Max-Q (8 GB, 140 W)
    \item NVIDIA H100 (PCIe)
    \item NVIDIA GH200
\end{itemize}


% ==========================================================
\section{Phase 1: CPU Baseline}

\subsection{Objective}
Establish a reproducible CPU reference for latency, throughput, and approximate GFLOPS before introducing CUDA-specific optimizations.

\subsection{Implementation}
The CPU code processes images in a loop and reports per-image statistics using structured timing and aggregate CSV output. This baseline is used as the control implementation for all speedup calculations in later phases.

\subsection{Experiments and Results}
The fixed benchmark set is executed multiple times to quantify variance. Reported numbers include average, median, and percentile latencies.

\begin{table}[H]
\centering
\begin{tabular}{lccc}
	oprule
Metric & Mean & Median & P95 \\
\midrule
Latency per image (ms) & TBD & TBD & TBD \\
Throughput (images/s) & TBD & - & - \\
GFLOPS (estimated) & TBD & - & - \\
\bottomrule
\end{tabular}
\caption{CPU baseline summary on the fixed benchmark image set.}
\end{table}

% ==========================================================
\section{Phase 2: Full CUDA Migration of the Pipeline}

\subsection{Objective}
Port all major stages of the processing chain to CUDA, not only preprocessing, so that the GPU remains active across the end-to-end pipeline minimizing host-device synchronization points.

\subsection{Scope}
The migration target includes:
\begin{itemize}
    \item Image preprocessing (grayscale/value extraction, contrast enhancement, blur, threshold)
    \item Plate localization stages
    \item Character extraction stages
\end{itemize}

% KNN remains a control path at this stage to isolate the impact of the core pipeline migration.

\subsection{Design Considerations}
The implementation emphasizes data layout, coalesced access, and minimizing host-device synchronization points. Intermediate results should remain on the device whenever possible.

\subsection{Experiments and Results}
Each migrated stage is validated against CPU output for correctness and compared by stage runtime.

\begin{table}[H]
\centering
\begin{tabular}{lcc}
	oprule
Stage & CPU Time (ms) & CUDA Time (ms) \\
\midrule
Preprocess & TBD & TBD \\
Plate detection & TBD & TBD \\
Character detection & TBD & TBD \\
\bottomrule
\end{tabular}
\caption{Stage-level runtime before and after CUDA migration.}
\end{table}


% ==========================================================
\section{Phase 3: Adaptive Batch Controller}

\subsection{Objective}
Reduce kernel idle time by dynamically adjusting batch size, keeping the GPU continuously supplied with data while avoiding transfer-side congestion.

\subsection{Controller Strategy}
The controller uses a global policy shared by all GPUs (extendable later to per-GPU tuning). Batch size can both increase and decrease based on observed transfer and compute behavior.

The controller parameters (growth thresholds, shrink thresholds, smoothing window, and min/max limits) are tuned experimentally.

\begin{lstlisting}[language=C++, caption={Adaptive batch-size controller (high-level logic)}]
int batchOfImages = 20;

while (has_more_images) {
    Stats s = run_batch(batchOfImages); // collect H2D, kernel, D2H, total times

    bool gpu_idle = s.kernel_idle_ratio > idle_threshold;
    bool transfer_limited = s.transfer_ratio > transfer_threshold;

    if (gpu_idle && !transfer_limited) {
        batchOfImages = min(batchOfImages + step_up, batch_max);
    } else if (transfer_limited || s.total_latency_ms > latency_target) {
        batchOfImages = max(batchOfImages - step_down, batch_min);
    }
}
\end{lstlisting}

\subsection{Experiments and Results}
Batch evolution is logged throughout each run to study convergence and stability.

\begin{table}[H]
\centering
\begin{tabular}{lccc}
	oprule
Device & Initial Batch & Final Mean Batch & Throughput Gain vs Static \\
\midrule
RTX 4070 Max-Q & 20 & TBD & TBD \\
H100 PCIe & 20 & TBD & TBD \\
GH200 & 20 & TBD & TBD \\
\bottomrule
\end{tabular}
\caption{Adaptive controller behavior and gain over static batching.}
\end{table}


% ==========================================================
\section{Phase 4: Asynchronous Pipeline (Single Stream to Multi Stream)}

\subsection{Objective}
Overlap data movement and compute to further reduce kernel idle periods and improve end-to-end throughput.

\subsection{Execution Path}
This phase proceeds in two steps:
\begin{itemize}
    \item Single-stream baseline: establish synchronized reference with CUDA timing events.
    \item Multi-stream execution: pipeline H2D transfer, kernel execution, and D2H copy across batches.
\end{itemize}

\subsection{Expected Effect}
The expected gain is highest on platforms where transfer overhead is significant. GH200 would show how CPU-GPU coupling changes transfer bottlenecks relative to PCIe devices.

\subsection{Experiments and Results}
\begin{table}[H]
\centering
\begin{tabular}{lccc}
	oprule
Device & Single Stream (img/s) & Multi Stream (img/s) & Improvement \\
\midrule
RTX 4070 Max-Q & TBD & TBD & TBD \\
H100 PCIe & TBD & TBD & TBD \\
GH200 & TBD & TBD & TBD \\
\bottomrule
\end{tabular}
\caption{Throughput impact of asynchronous multi-stream execution.}
\end{table}


% ==========================================================
% \section{Phase 5: KNN Control and Experiments}

% \subsection{Objective}
% Use KNN as a controlled reference path, then explore optional optimizations while tracking both accuracy and performance impact.

% \subsection{Plan}
% Phase 5 is split into two tracks:
% \begin{itemize}
%     \item Control track: preserve original KNN behavior for direct comparability with earlier phases.
%     \item Experimental track: test selective optimizations (data layout, batching of descriptor comparisons, and optional CUDA-assisted variants).
% \end{itemize}

% \subsection{Evaluation}
% KNN experiments are accepted only if they preserve detection quality on the fixed benchmark set and improve either latency or throughput.

% \begin{table}[H]
% \centering
% \begin{tabular}{lccc}
% 	oprule
% KNN Variant & Accuracy Metric & Time per Plate (ms) & Relative Speedup \\
% \midrule
% Control (reference) & TBD & TBD & 1.00x \\
% Experiment A & TBD & TBD & TBD \\
% Experiment B & TBD & TBD & TBD \\
% \bottomrule
% \end{tabular}
% \caption{KNN control and experiment results.}
% \end{table}


% ==========================================================
\section{Cross-Device Comparison Summary}

\subsection{Reporting Format}
Final comparison is presented for the fixed benchmark set across all platforms, using the same controller logic and the same evaluation protocol.

\begin{table}[H]
\centering
\begin{tabular}{lcccc}	
\toprule
Platform & Latency (ms/img) & Throughput (img/s) & End-to-End Speedup vs CPU & Notes \\
\midrule
CPU baseline & TBD & TBD & 1.00x & reference \\
RTX 4070 Max-Q & TBD & TBD & TBD & mobile GPU \\
H100 PCIe & TBD & TBD & TBD & datacenter PCIe \\
GH200 & TBD & TBD & TBD & CPU-GPU integrated platform \\
\bottomrule
\end{tabular}
\caption{Final platform comparison on a fixed image set.}
\end{table}

\subsection{Conclusions}
This phased methodology allows clear attribution of gains: baseline quality and timing from CPU, acceleration from full CUDA migration, 
utilization improvements from adaptive batching, overlap benefits from asynchronous execution. 
% and controlled classifier analysis in the final KNN phase.

\subsection*{Addendum: Individual responsibilities in this assignment}

\begin{table}[H]
\centering
\begin{tabular}{|c|c|c|}
\hline
	extbf{Part} & \textbf{Student No.} & \textbf{Name} \\
\hline
I-V & s256878 & Dragos-Daniel Bonaparte \\
\hline
\end{tabular}
\end{table}
